\section{관련 연구}
\label{sec:관련연구}

\subsection{FANET과 지리적 라우팅}
FANET 환경에서의 라우팅 프로토콜은 UAV의 3차원 고속 이동성과 채널 페이딩, 물리적 통신 도달 범위 제약에 매우 실시간으로 대응해야 한다. 전통적인 토폴로지 기반 라우팅과 달리 GPSR~\cite{gpsr}과 같은 지리적 라우팅 프로토콜은 이웃 무인기들의 공간 위치만을 활용하여 차홉을 그리디하게 지정한다. 제어 오버헤드가 없고 즉각적인 전송 결정을 내릴 수 있어 FANET의 강력한 기준선 프로토콜로 사용되나, 라우팅 홀 지역이나 채널 감쇠, 고이동성 링크 끊김 상황에서는 패킷이 손실되는 근본적인 한계를 지닌다.

\subsection{예측적 지리적 라우팅}
지리적 라우팅의 한계를 보완하기 위해 노드의 속도 벡터, 가속도 등을 활용하여 미래의 링크 지속 시간(Link Expiration Time, LET)이나 토폴로지의 변화를 미리 예측하는 예측적 지리적 라우팅 기법들이 제시되었다. 특히 Evo-QGeo~\cite{evoqgeo} 등은 진화 연산이나 기하학적 예측 모형을 결합해 미래 시점의 링크 수명과 경로 상의 안정성을 종합 평가하여 라우팅 홀을 능동적으로 우회한다. 이 방식은 PDR을 대폭 끌어올리지만, 통신 패킷 내에 대량의 이동성 파라미터를 추가하여 주변 이웃들에게 주기적으로 공유해야 하는 제어 시그널 오버헤드를 유발한다.

\subsection{강화학습 및 다중 에이전트 강화학습 라우팅}
강화학습을 라우팅에 적용할 경우 채널 에러율, 대기열 지연, 홉 수 등 비선형적인 다중 메트릭을 수렴된 Q-값이나 정책 네트워크를 통해 최적화할 수 있다. Q-Routing은 실시간 지연 시간을 반영해 최단 경로를 우회하는 이점이 있으며, DRAMA~\cite{drama}와 같은 기법은 다중 에이전트 환경에서 에이전트 간의 '창발적 통신(Emergent Communication)' 학습을 도입해 주변 상황을 서로 교환하여 패킷 전송을 협력적으로 제어한다. 그러나 이러한 방식 역시 배포 단계에서 무선 대역폭을 점유하는 실시간 제어 메시지 전송이 지속되어야 하므로 실질적인 UAV swarms 시스템 배포에는 많은 제약이 따른다.

\subsection{GNN 기반 라우팅 및 CTDE}
네트워크 토폴로지는 전형적인 그래프 구조이므로, Graph Neural Networks (GNN)을 활용해 링크 간 가중치나 통신 특징 공간을 임베딩하는 연구들이 활발히 진행되어 왔다. 중앙 집중식 학습 및 분산 실행(CTDE) 구조 하에서 GLo-MAPPO~\cite{glo_mappo} 등은 오프라인 학습 시점에 시뮬레이터 내에서 전역 그래프 토폴로지 정보를 입력하여 최적의 제어 정책을 학습하고, 실제 가동 시에는 개별 분산 에이전트의 로컬 상태만을 입력하여 구동되도록 설계된다. 그러나 분산 UAV 시스템 내부에서 매 홉마다 딥러닝 기반의 GNN 추론 연산을 즉시 수행하는 것은 심각한 계산 지연을 야기한다. 본 연구는 정책 증류 기법을 활용해 전역 정보 기반 GNN 교사의 출력을 로컬 1-hop 피처 입력만을 사용하는 가벼운 MLP 학생 네트워크로 변환함으로써 연산 지연을 줄이고 실시간성을 보장한다.
